%【紙面設定】
%b4,横書き,2段
\documentclass[paper=b4j,landscape,twocolumn,fleqn]{jlreq}
\special{papersize=\the\paperwidth,\the\paperheight}
\setlength{\columnseprule}{0.4pt}
\pagestyle{empty}
\usepackage[margin=10truemm]{geometry}
\usepackage[dvipdfmx]{graphicx}
\usepackage{amsmath}
\ModifyHeading{section}{lines=1}

%【本文】
\begin{document}
\section{五則演算 解答}
\begin{align*}
5+8&=13 & 9-4&=5 & 2*3&=6 & 6/2&=3 & 5\%2&=1\\
-9+4&=-5 & 3-10&=-7 & -4*5&=-20 & -10/5&=-2 & 10\%3&=1\\
15+27&=42 & 50-18&=32 & 12*4&=48 & 48/4&=12 & 25\%4&=1\\
-50+12&=-38 & 20-85&=-65 & 15*-3&=-45 & 100/-25&=-4 & 100\%7&=2\\
1.2+3.8&=5 & 4.5-1.2&=3.3 & 0.5*6&=3 & 4.8/2&=2.4 & 50\%6&=2\\
-4.5+1.2&=-3.3 & 0.8-3.5&=-2.7 & -1.2*4&=-4.8 & -7.2/0.9&=-8 & 80\%9&=8\\
125+380&=505 & 500-125&=375 & 25*12&=300 & 144/12&=12 & 123\%10&=3\\
-600+250&=-350 & 100-456&=-356 & 14*-10&=-140 & -117/9&=-13 & 250\%15&=10\\
0.45+0.55&=1 & 1.25-0.75&=0.5 & 0.8*0.5&=0.4 & 0.75/0.25&=3 & 1000\%3&=1\\
-12.5+7.5&=-5 & 1000-1500&=-500 & -15*1.2&=-18 & 102/-3&=-34 & 1024\%10&=4
\end{align*}

\section{ルートと乗数の計算 解答}
\begin{align*}
2^2&=4 & 2^2*2^1&=8 & \sqrt{4}&=2 & \sqrt{4}*\sqrt{9}&=6\\
3^2&=9 & 3^3/3^1&=9 & \sqrt{9}&=3 & \sqrt{2}*\sqrt{8}&=4\\
2^4&=16 & 2^2*2^2&=16 & \sqrt{16}&=4 & \sqrt{12}/\sqrt{3}&=2\\
5^2&=25 & 10^2/10^1&=10 & \sqrt{25}&=5 & \sqrt{3}*\sqrt{12}&=6\\
4^3&=64 & 2^3*2^2&=32 & \sqrt{49}&=7 & \sqrt{18}/\sqrt{2}&=3\\
6^2&=36 & 5^3/5^2&=5 & \sqrt{64}&=8 & \sqrt{2}*\sqrt{32}&=8\\
10^2&=100 & 2^2*2^4&=64 & \sqrt{100}&=10 & \sqrt{50}/\sqrt{2}&=5\\
12^2&=144 & 3^4/3^2&=9 & \sqrt{144}&=12 & \sqrt{72}/\sqrt{2}&=6
\end{align*}

\section{分数の計算 解答}
\begin{align*}
\frac{1}{4}+\frac{1}{4}&=\frac{1}{2} & \frac{4}{5}-\frac{1}{5}&=\frac{3}{5} & \frac{1}{3}\times\frac{1}{2}&=\frac{1}{6} & \frac{2}{5}\div2&=\frac{1}{5}\\
-\frac{2}{3}+\frac{1}{3}&=-\frac{1}{3} & \frac{1}{4}-\frac{3}{4}&=-\frac{1}{2} & -\frac{1}{2}\times\frac{1}{4}&=-\frac{1}{8} & \frac{1}{5}\div(-2)&=-\frac{1}{10}\\
\frac{1}{2}+\frac{1}{3}&=\frac{5}{6} & \frac{1}{2}-\frac{1}{4}&=\frac{1}{4} & \frac{2}{3}\times\frac{3}{5}&=\frac{2}{5} & \frac{3}{4}\div2&=\frac{3}{8}\\
-\frac{1}{2}+\frac{1}{6}&=-\frac{1}{3} & \frac{1}{6}-\frac{1}{2}&=-\frac{1}{3} & \frac{2}{5}\times\left(-\frac{1}{4}\right)&=-\frac{1}{10} & -\frac{3}{4}\div3&=-\frac{1}{4}\\
\frac{3}{4}+\frac{1}{6}&=\frac{11}{12} & \frac{5}{6}-\frac{1}{4}&=\frac{7}{12} & \frac{5}{8}\times\frac{4}{15}&=\frac{1}{6} & \frac{2}{3}\div\frac{1}{2}&=\frac{4}{3}\\
-\frac{1}{2}-\frac{1}{3}-\frac{1}{6}&=-1 & \frac{1}{10}-\frac{4}{15}&=-\frac{1}{6} & -\frac{15}{16}\times\frac{4}{5}&=-\frac{3}{4} & 10\div\left(-\frac{5}{2}\right)&=-4
\end{align*}

\newpage
\section{方程式の計算 解答}
\begin{align*}
x+5&=8\Rightarrow x=3 & 3x&=12\Rightarrow x=4 & x+10&=4\Rightarrow x=-6\\
2x-3&=7\Rightarrow x=5 & 15-x&=20\Rightarrow x=-5 & 4x+16&=0\Rightarrow x=-4\\
3x+4&=x+10\Rightarrow x=3 & 2x-5&=5x+7\Rightarrow x=-4 & 4x-8&=6x+2\Rightarrow x=-5\\
2(x-3)&=4\Rightarrow x=5 & 3(x+4)&=-6\Rightarrow x=-6 & 5(x+1)&=2x-4\Rightarrow x=-3\\
0.2x&=4\Rightarrow x=20 & \frac{x}{3}-2&=-4\Rightarrow x=-6 & 0.5x+3&=1\Rightarrow x=-4
\end{align*}

\section{式の展開 解答}
\begin{align*}
5(x+3)&=5x+15 & (x+2)(x+6)&=x^2+8x+12\\
-2(4x-7)&=-8x+14 & (x-3)(x+5)&=x^2+2x-15\\
(x+8)^2&=x^2+16x+64 & (x-9)^2&=x^2-18x+81\\
(x+12)(x-12)&=x^2-144 & (2x+3)(x+2)&=2x^2+7x+6\\
(3x-1)(2x+5)&=6x^2+13x-5 & (4x-3)^2&=16x^2-24x+9
\end{align*}

\section{因数分解 解答}
\begin{align*}
2x+10&=2(x+5) & x^2+7x+10&=(x+5)(x+2)\\
x^2-2x-8&=(x-4)(x+2) & x^2-5x+4&=(x-1)(x-4)\\
x^2+6x+9&=(x+3)^2 & x^2-25&=(x-5)(x+5)\\
x^2+x-20&=(x+5)(x-4) & x^2+10x+25&=(x+5)^2\\
x^2-81&=(x-9)(x+9) & x^2+13x+42&=(x+6)(x+7)
\end{align*}

\section{文字と式 解答}
\begin{align*}
&1.\ 5a+3b=c\\
&2.\ 1000-3x=y\\
&3.\ vt=d\\
&4.\ 100-xt=y\\
&5.\ 1.1(2x+5y)=z\\
&6.\ \frac{a+b+c+d}{4}=m\\
&7.\ x-vt=p\\
&8.\ 10a\left(1-\frac{x}{100}\right)=b
\end{align*}
\end{document}
